% Slajd 10 – od prototypu do usługi: co jest, co dalej, warunki kolejnego miasta (pkt 6 wyzwania).
\begin{frame}{Od prototypu do usługi}
  \begin{columns}[T,onlytextwidth]
    \begin{column}{0.45\textwidth}
      \begin{exampleblock}{Mamy dziś}
        \footnotesize
        \begin{itemize}
          \item Prototyp na telefon i~komputer: 6\,545 miejsc, 27 atrybutów, 11 warstw, mapa barier, asystent, 4~języki.
          \item Panel właściciela i~panel administratora (kolejka, statystyki, braki danych).
          \item Źródło, data i~status przy każdym fakcie; brak danych pokazany jako brak.
          \item Sprawdzone kryteria WCAG 2.2 A/AA w~głównym scenariuszu; 503 testy, Docker, Ansible, HTTPS.
        \end{itemize}
      \end{exampleblock}
    \end{column}
    \begin{column}{0.53\textwidth}
      \begin{block}{Plan na 3--6 miesięcy}
        \footnotesize
        \begin{enumerate}
          \item \textbf{Pilotaż} z~organizacją osób z~niepełnosprawnościami, 2--3 hotelami i~instytucjami kultury;
                audyt z~użytkownikami.
          \item \textbf{Produkcyjne logowanie} (poczta, opcjonalnie OAuth), powiadomienia, tryb offline.
          \item \textbf{AI do weryfikacji zdjęć i~tekstu} -- propozycje do potwierdzenia, nigdy fakty; silnik zaufania z~Rampy.
          \item \textbf{Open API i~widget} dla właścicieli; integracja z~systemami rezerwacyjnymi.
          \item \textbf{Drugie miasto:} plik miasta + wyciąg OSM + warstwy miejskie; warunek: otwarte dane (lub OSM)
                i~lokalny partner do moderacji.
        \end{enumerate}
      \end{block}
    \end{column}
  \end{columns}
  \vspace{1.5mm}
  \centering
  {\bfseries\color{accNavy}Nie kolejna statyczna mapa -- żywy, weryfikowalny obraz dostępności miasta.\par}
  \vspace{0.8mm}
  \muted{Demo: \DemoUrl\ \textperiodcentered\ Kod: \RepoUrl\ \textperiodcentered\ Film: \VideoUrl}

  \note[item]{Zakończyć zdaniem z~tego slajdu i~zaprosić do stolika z~telefonem – niech jury samo kliknie „Zgłoś problem”.}
  \note[item]{Pytanie o~zespół po hackathonie: operator (spółka / fundacja), moderacja w~panelu, finansowanie z~planu Pro i~wdrożeń white-label.}
  \note[item]{Pytanie o~drugie miasto: potrzebne są granice i~nazwa (plik miasta), wyciąg OSM z~Geofabrik oraz – jeśli miasto je ma – otwarte warstwy (przystanki, toalety, parkingi). Bez danych miejskich działa sam OSM.}
\end{frame}
